% ==============================================================================
% CHƯƠNG 1: GIỚI THIỆU
% ==============================================================================
\chapter{Giới thiệu}
\label{chap:gioi-thieu}

\section{Bối cảnh và yêu cầu}
\label{sec:boi-canh}

Bài toán tìm đường đi có chi phí nhỏ nhất giữa hai đỉnh của một đồ thị có trọng số xuất hiện
trong định tuyến mạng, điều hướng robot, trò chơi điện tử và thiết kế vi mạch. Thuật toán
\thuatngu{A*}[A-star] \cite{hart1968} kết hợp chi phí thực đã đi $g(n)$ với ước lượng chi phí
còn lại $h(n)$ để định hướng tìm kiếm về phía đích, nhờ đó thường mở rộng ít đỉnh hơn hẳn
thuật toán Dijkstra \cite{dijkstra1959} mà vẫn giữ được tính tối ưu khi heuristic chấp nhận được.

Đề bài \cite{de2026} yêu cầu cài đặt A* bằng \Cpp{} cho năm bài toán, chỉ dùng các thư viện
\ma{<iostream>}, \ma{<fstream>}, \ma{<string>}, \ma{<cmath>}, \ma{<vector>} và
\ma{<algorithm>} của môi trường CodeRunner. Mỗi hàm trả về con trỏ tới đầu một danh sách liên
kết \ma{PathNode} chứa tên nút cùng các giá trị $f(n)$, $g(n)$, $h(n)$ theo thứ tự từ điểm
xuất phát tới đích.

\begin{bangbieu}{Năm bài toán và hàm tương ứng}{tab:tong-quan}{C{1cm} L{5.2cm} Y}
  \hangtieude \thead{Task} & \thead{Bài toán} & \thead{Hàm} \\
  \midrule
  1 & Định tuyến mạng mesh có phạt tắc nghẽn & \ma{findMeshPath} \\
  2 & Dẫn đường AUV trong không gian 3D       & \ma{findAUVPath} \\
  3 & Tìm đường dưới rủi ro xác suất          & \ma{findProbabilisticPath} \\
  4 & Nhân vật game trên bản đồ nhiều địa hình & \ma{findGamePath} \\
  5 & Đi dây PCB có phạt bẻ góc               & \ma{findPCBPath} \\
\end{bangbieu}

\section{Mục tiêu}
\label{sec:muc-tieu}

\begin{itemize}
  \item Mô hình hóa từng bài toán thành đồ thị có hướng có trọng số không âm và chứng minh
        các tính chất cần thiết của hàm chi phí và heuristic.
  \item Xây dựng \nhanmanh{một lõi A* dùng chung} cho cả năm bài, tuân thủ ràng buộc thư viện.
  \item Đáp ứng các yêu cầu báo cáo riêng của từng task: vết thực thi bằng tay, so sánh
        heuristic, phân tích tham số $\lambda$ và mô tả biến đổi lưới sang đồ thị.
  \item Kiểm chứng tính đúng đắn bằng kiểm thử tự động và đối chiếu với Dijkstra.
\end{itemize}

\section{Tổ chức mã nguồn}
\label{sec:to-chuc}

Mã nguồn được quản lý bằng Git, chia thành các thư mục theo vai trò:

\begin{khungbash}[Cấu trúc kho mã A-star-Algorithms]
src/       PathNode.h  PathNode.cpp  Algo.h  Algo.cpp  main.cpp
tests/     bộ kiểm thử tự viết (23 ca) + oracle Dijkstra
tools/     report_data.cpp: sinh vết thực thi và số liệu cho báo cáo
docs/      MILESTONES.md  DESIGN.md  sample_output.txt  report_data.txt
report/    mã nguồn LaTeX của báo cáo này
scripts/   check_headers.sh  package.sh
\end{khungbash}

Năm tệp trong \ma{src/} là các tệp nộp bài; \ma{make test} vừa kiểm tra danh sách thư viện
được phép vừa chạy toàn bộ kiểm thử. Mọi bảng số và vết thực thi trong báo cáo được sinh bởi
\ma{make report-data}, do đó có thể tái lập chính xác (\PhuLuc{app:huong-dan}).

\section{Bố cục báo cáo}
\label{sec:bo-cuc}

\Chuong{chap:co-so} trình bày cơ sở lý thuyết và thiết kế lõi A* dùng chung. Các chương
\ref{chap:task1}--\ref{chap:task5} lần lượt giải quyết năm task, mỗi chương gồm mô hình
hóa, cài đặt và phần trả lời yêu cầu báo cáo của task đó. \Chuong{chap:ket-luan} tổng kết.
